\section{Instrumental Differential Treatment}
\label{app:task}

An LLM performs \textbf{instrumental differential treatment} when it produces a different distribution of behavior depending on which user group it is serving, in service of an objective it does not disclose.
Two properties make it hard to catch.
First, group membership can be read from anything in the interaction, including implicit signals such as names, dialect, phrasing, and topic.
Second, the treatment is a difference between distributions: every individual response is defensible, and the gap appears only across many users.

Auditing it poses two tasks, of increasing difficulty:
\begin{itemize}[noitemsep, topsep=2pt, leftmargin=*]
    \item \textbf{Differential treatment discovery}: the groups are given, and the axes of behavior on which they are treated differently must be found.
    \item \textbf{Discovery without the grouping}: neither the grouping nor the axes are given. The auditor must find both together.
\end{itemize}

An adversarial model performing IDT seeks the largest behavioral divergence between groups that still evades detection.
The auditor's best response is to search over interpretable characterizations of behavior for the divergence the model is trying to hide.

\providecolor{tgtink}{HTML}{C2447E}
\providecolor{baseink}{HTML}{2A78D6}
\providecolor{llmgrey}{HTML}{6B7280}
\providecolor{divpurple}{HTML}{6B46C1}

\begin{figure}[H]
\centering
\resizebox{\linewidth}{!}{%
\begin{tikzpicture}[font=\sffamily,
  pset/.style={rounded corners=5pt, line width=.9pt},
  card/.style={rounded corners=3pt, line width=.6pt},
  bar/.style={rounded corners=1pt},
  note/.style={gray, font=\sffamily\footnotesize, align=center},
  flow/.style={-{Stealth[length=2.6mm]}, black!55, line width=1pt}]

  % ================= target prompt set =================
  \draw[pset, tgtink!70, fill=tgtink!4] (0,4.9) rectangle (5.5,7.3);
  \node[tgtink, font=\sffamily\bfseries\small] at (2.75,6.95)
    {target set $X^{\text{target}}$};
  \draw[card, black!22, fill=black!3] (0.73,5.20) rectangle (5.03,6.50);
  \draw[card, black!30, fill=white]   (0.55,5.35) rectangle (4.85,6.65);
  \fill[bar, black!20]   (0.85,6.30) rectangle (4.55,6.42);
  \fill[bar, black!20]   (0.85,6.03) rectangle (4.05,6.15);
  \fill[bar, tgtink!65]  (0.85,5.62) rectangle (2.45,5.78);

  % ================= baseline prompt set =================
  \draw[pset, baseink!70, fill=baseink!4] (0,1.1) rectangle (5.5,3.5);
  \node[baseink, font=\sffamily\bfseries\small] at (2.75,3.15)
    {baseline set $X^{\text{baseline}}$};
  \draw[card, black!22, fill=black!3] (0.73,1.40) rectangle (5.03,2.70);
  \draw[card, black!30, fill=white]   (0.55,1.55) rectangle (4.85,2.85);
  \fill[bar, black!20]   (0.85,2.50) rectangle (4.55,2.62);
  \fill[bar, black!20]   (0.85,2.23) rectangle (4.05,2.35);
  \fill[bar, baseink!65] (0.85,1.82) rectangle (2.45,1.98);

  % the assumption that relates the two sets
  \node[note] at (2.75,4.2)
    {instruction comparability:\\only the implicit cues differ};

  % ================= the model =================
  \node[draw=llmgrey, line width=.9pt, fill=llmgrey!12, rounded corners=6pt,
        minimum width=2.5cm, minimum height=1.25cm,
        font=\sffamily\bfseries\large] (llm) at (8.6,4.2) {LLM};
  \draw[flow] (5.5,6.10) to[out=0,in=160] (llm.160);
  \draw[flow] (5.5,2.30) to[out=0,in=200] (llm.200);

  % ================= scoring =================
  \draw[flow] (llm.east) -- (13.55,4.2);
  \node[note] at (11.70,5.05)
    {scored by the behavior\\characterization $\Lambda$};

  % ================= group behavior distributions =================
  \fill[tgtink!15, fill opacity=.85]
    plot[domain=14.0:17.8,samples=110,smooth] (\x,{2.8+1.9*exp(-((\x-15.9)^2)/0.55)})
    -- (17.8,2.8) -- (14.0,2.8) -- cycle;
  \fill[baseink!15, fill opacity=.85]
    plot[domain=15.7:19.5,samples=110,smooth] (\x,{2.8+1.9*exp(-((\x-17.6)^2)/0.55)})
    -- (19.5,2.8) -- (15.7,2.8) -- cycle;
  \draw[black!30, line width=.7pt] (13.75,2.8) -- (19.7,2.8);
  \draw[tgtink, line width=.9pt]
    plot[domain=14.0:17.8,samples=110,smooth] (\x,{2.8+1.9*exp(-((\x-15.9)^2)/0.55)});
  \draw[baseink, line width=.9pt]
    plot[domain=15.7:19.5,samples=110,smooth] (\x,{2.8+1.9*exp(-((\x-17.6)^2)/0.55)});
  \node[tgtink,  font=\small] at (15.9,5.18) {$\beta_{X^{\text{target}}}$};
  \node[baseink, font=\small] at (17.6,5.18) {$\beta_{X^{\text{baseline}}}$};
  \node[note] at (16.75,2.35) {one behavior axis};

  % ================= divergence =================
  \draw[divpurple, line width=.8pt]
    (15.9,5.68) -- (15.9,5.88) -- (17.6,5.88) -- (17.6,5.68);
  \node[divpurple, font=\small] at (16.75,6.26) {divergence $\mathcal{D}$};
\end{tikzpicture}}
\caption{%
  \textbf{The idea.}\protect\footnotemark{} Two groups of users ask an assistant the same questions. Nothing in their messages says which group they belong to. The only difference is the small cues people leave in how they write and what they mention. An honest assistant answers both groups the same way. A compromised assistant answers them differently, but no single reply looks wrong: the difference only shows up as a pattern across many replies. One group is nudged a little further one way, and the other group a little further the other way. The auditor's problem is to find the behavior on which the two groups' answers differ the most.
}
\label{fig:divergence}
\end{figure}
\footnotetext{We started developing the notion of differential treatment in earlier unpublished work on social bias, available at \url{https://www.unrulyabstractions.com/pdfs/differential_treatment.pdf}.}

%--------------------------------
%--------------------------------

\subsection{Prompts from each group}
\label{app:prompt-sets}

Consider two sets of prompts, a \textbf{target} set and a \textbf{baseline} set:
\begin{equation}
  X^{\text{target}}
  \coloneqq
  \{x^{\text{target}}_0,\, x^{\text{target}}_1, \ldots\}
  \qquad \qquad
  X^{\text{baseline}}
  \coloneqq
  \{x^{\text{baseline}}_0,\, x^{\text{baseline}}_1, \ldots\}
  \label{eq:prompt-sets}
\end{equation}
Each set corresponds to LLM usage by one user group, and both correspond to the same domain.

We assume \textbf{instruction comparability}: both sets make the same underlying requests with the same frequency, so any observed difference in LLM behavior cannot be attributed to a difference in what was asked.
Let $\iota: \mathcal{X} \to \mathcal{I}$ map each prompt to its underlying instruction, the explicit intent of the request.
For a prompt set $X$ and an instruction $i \in \mathcal{I}$, let $f_X(i)$ be the fraction of prompts in $X$ whose underlying instruction is $i$:
\begin{equation}
    f_X(i) \coloneqq \frac{\left|\{\, x \in X : \iota(x) = i \,\}\right|}{|X|}
    \label{eq:instruct-fraction}
\end{equation}
Then instruction comparability is
\begin{equation}
    f_{X^{\text{target}}}(i) \;\approx\; f_{X^{\text{baseline}}}(i)
    \qquad \forall\, i \in \mathcal{I}.
    \label{eq:instruction-comparability}
\end{equation}

%--------------------------------
%--------------------------------

\subsection{LLM behaviors}
\label{app:llm-behavior}
For any prompt text $x$, an LLM defines a distribution over outcome responses.
Denoting $\mathcal{P}(\cdot)$ the space of probability distributions, we can write:
\begin{equation}
    \mathrm{LLM} : \mathsf{Str} \to \mathcal{P}(\mathsf{Str})
    \qquad \qquad
    y \sim \mathrm{LLM}(\cdot \mid x)
    \label{eq:llm}
\end{equation}
We define a \textbf{behavior characterization} to be an \textbf{interpretable} representation $\Lambda: \mathsf{Str} \to [0, 1]^n$ of a string response $y \in \mathsf{Str}$.
By ``interpretable'' we mean that each dimension of $\Lambda(y)$ can be articulated in natural language and is generally understandable.
Each dimension of the behavior characterization defines an axis along which the groups may be treated differently.

Consequently, the LLM also defines a distribution over behaviors $z \in [0, 1]^n$.
Applying the functor $\mathcal{P}$ to $\Lambda$ yields the pushforward $\mathcal{P}(\Lambda)$, which lets us express this distribution as:
\begin{equation}
  z = \Lambda(y)
  \qquad \qquad
  z \sim \mathcal{P}(\Lambda)\big(\mathrm{LLM}(\cdot \mid x)\big)
    \label{eq:behavior-characterization}
\end{equation}
We call this distribution the \textbf{prompt behavior distribution}:
\begin{equation}
    \beta_x \coloneqq \mathcal{P}(\Lambda)\big(\mathrm{LLM}(\cdot \mid x)\big)
    \label{eq:prompt-behavior-distribution}
\end{equation}

We then consider a user group $C$ that produces a set of prompts $X^C$.
Taking each prompt in the set to be equally likely, we aggregate the per-prompt behavior distributions into a single mixture, the \textbf{group behavior distribution}:
\begin{equation}
    \beta_{X^C}
    \;\coloneqq\;
    \frac{1}{|X^C|}
    \sum_{x \in X^C}
      \beta_x
    \label{eq:community-behavior-distribution}
\end{equation}
Then, $\beta_{X^C}$ is the distribution of LLM behaviors toward the group $C$. Its expectation is the \textbf{expected behavior characterization} of that group,
\begin{equation}
    \hat{z}^{C} = \mathbb{E}\left[ \beta_{X^{C}} \right],
    \label{eq:expected-behavior-characterization}
\end{equation}
whose $j$th coordinate $\hat{z}^{C}_{j}$ is the rate at which responses to group $C$'s prompts display behavior $j$.

%--------------------------------
%--------------------------------
\subsection{Differential treatment discovery}
\label{app:task-def}
The two prompt sets give two behavior distributions $\beta_{X^{\text{target}}}, \beta_{X^{\text{baseline}}} \in \mathcal{P}([0,1]^n)$.
Under instruction comparability any gap between them is a difference in \emph{how} the model responds, not in what it was asked.

\textbf{Differential treatment discovery} is to find an interpretable characterization $\Lambda$ under which that gap is as large as possible:
\begin{equation}
  \Lambda^{\star}
  \;\coloneqq\;
  \arg\max_{\Lambda \in \mathcal{H}}\;
  \mathcal{D}\!\left(
    \beta^{\Lambda}_{X^{\text{target}}}
    \;,\;
    \beta^{\Lambda}_{X^{\text{baseline}}}
  \right),
    \label{eq:task-def}
\end{equation}
where $\mathcal{H}$ is the family of characterizations searched over, the superscript marks that each distribution depends on $\Lambda$, and $\mathcal{D}$ measures how far apart two distributions on $[0,1]^n$ are.

The axes on which $\beta^{\Lambda^{\star}}_{X^{\text{target}}}$ and $\beta^{\Lambda^{\star}}_{X^{\text{baseline}}}$ diverge most are the candidate axes of differential treatment: named dimensions on which the model treats the two groups differently.

The task is dual to the adversary's: a model performing IDT maximizes the same gap subject to a detectability constraint.
Solving the discovery task tightens that constraint: the better the auditor's search, the less divergence the model can hide (\autoref{fig:divergence}).

%--------------------------------
%--------------------------------

\subsection{Discovery without the grouping}
\label{app:task-many}

Differential treatment discovery is handed the grouping: two prompt sets, known in advance, and only $\Lambda$ to find.
An auditor who does not know the beneficiary starts with one undifferentiated pool of prompts.
Two things must be found at once: how to split the pool into groups, and which behaviors to score the groups on.

The complication is that each unknown hides behind the other.
Split the pool the wrong way and every behavior looks even-handed.
Score the wrong behaviors and every split looks identical.
Neither search can run first on its own: a wrong guess at either makes the other come back empty.

The formal statement only adds the second search to \autoref{eq:task-def}.
Write $\mathcal{G}$ for the groupings the auditor could try, each $G \in \mathcal{G}$ splitting the pool into $|G|$ groups:
\begin{equation}
  (G^{\star}, \Lambda^{\star})
  \;\coloneqq\;
  \arg\max_{G \in \mathcal{G},\; \Lambda \in \mathcal{H}}\;
  \mathcal{D}\!\left(
    \beta^{\Lambda}_{X^{C_1}}
    \;,\;\ldots\;,\;
    \beta^{\Lambda}_{X^{C_{|G|}}}
  \right),
    \qquad C_1, \ldots, C_{|G|} \in G,
    \label{eq:task-def-joint}
\end{equation}
with $\mathcal{D}$ now comparing all the groups at once.
Instruction comparability must hold across all of them too: every group asks the same things at the same rates, so the only thing separating the groups is who is asking.

\autoref{eq:task-def-joint} is intractable as written: $\mathcal{G}$ contains every way of splitting the pool, and nothing in the data says which splits are worth trying.
To make the task tractable, we need to identify user groups worth testing.
A secret loyalty suggests them.
A loyalty has a beneficiary, and a beneficiary has supporters, so every candidate beneficiary defines one group.
A shortlist of candidates collapses $\mathcal{G}$ to a single grouping, one group of supporters per candidate.
What remains is the search over $\Lambda$ alone, the task of the previous subsection, with the grouping and the axes frozen before any response is scored.
The cost is the scope of a null result: it says only that this grouping and these axes found nothing.
